\documentclass[11pt,a4paper]{article}
\usepackage[T1]{fontenc}
\usepackage[utf8]{inputenc}
\usepackage{lmodern}
\usepackage{amsmath,amssymb}
\usepackage[margin=25mm]{geometry}
\usepackage[hidelinks]{hyperref}
\hypersetup{pdftitle={Constructing a 0+1 field theory: the quartic oscillator and the Yang--Mills matrix model},pdfauthor={navstokgap research working notes}}
\newcommand{\tightlist}{\setlength{\itemsep}{0pt}\setlength{\parskip}{0pt}}
\setlength{\emergencystretch}{3em}
\setcounter{secnumdepth}{0}
\begin{document}
\hypertarget{constructing-a-01-field-theory-the-quartic-oscillator-and-the-yangmills-matrix-model}{%
\section{Constructing a 0+1 field theory: the quartic oscillator and the
Yang--Mills matrix
model}\label{constructing-a-01-field-theory-the-quartic-oscillator-and-the-yangmills-matrix-model}}

\textbf{Result, 2026-10-27; standard construction assembled here, not a
novel existence theorem.} For any supplied \(m,\lambda,\hbar>0\), the
one-component \(\phi^4_1\) theory has a continuous, stationary Euclidean
field on the full time line, a reflection-positive bounded-cylinder
algebra, a non-Gaussian vacuum, and a reconstructed positive-energy
quantum system with a unique vacuum and a positive spectral gap. The
finite-time lattice measures with ground-state boundary weights converge
under arbitrary time refinement. Their real-time counterpart is the
nonrelativistic quartic oscillator. The same construction applies to the
existing \(SU(2)\) bosonic matrix Hamiltonian, including its
gauge-singlet sector, using the compact-resolvent result already proved
in \href{low-dimensional-mass-gap.md}{low-dimensional mass gap, Theorem
6}.

The one-dimensional Euclidean field has only time as a base-space
coordinate, so spacelike locality and Lorentz boosts impose no
nontrivial condition. A relativistic particle may instead be described
by a one-dimensional worldline theory with a spatial target and a
different Hamiltonian. Neither construction derives \(\hbar>0\), proves
the original 1998 consistency conjecture, or constructs four-dimensional
Yang--Mills. The statements below use established Feynman--Kac,
semigroup and spectral theorems; the assembly is written here, without a
separate referee check.

\hypertarget{model-and-prior-work}{%
\subsection{1. Model and prior work}\label{model-and-prior-work}}

Put \[
 V(q)=\frac{\lambda}{4}q^4,\qquad
 S_E[q]=\int\left(\frac m2\dot q^{\,2}+V(q)\right)d\tau.
 \tag{1}
\] The symbolic \(e^{-S_E/\hbar}\mathcal Dq\) below means an actual
Brownian-bridge tilt or its stationary extension, not a Lebesgue measure
on path space. Brownian paths are almost surely nowhere differentiable;
the kinetic integral in (1) is therefore a mnemonic for the reference
measure, not an ordinary pathwise integral.

The \href{rivero-1998-conjecture-central-forces.md}{Rivero 1998 test},
Theorems 2(b) and 3, already proves the fixed-\(\hbar\) real-time
product limit and the zero-endpoint Euclidean bridge limit for the
potential \(\lambda_{\rm old}q^4\); here
\(\lambda_{\rm old}=\lambda/4\). Its zero-endpoint bound also permits
\(\hbar\downarrow0\), so it does not select a quantum action unit. The
\href{newton-nonlinear-time-cell.md}{nonlinear cell} is one exact
elimination step for (1), with its full normalized interaction retained.
Its one-step estimates alone were not a stationary continuum
construction.

Feynman--Kac and the \(P(\phi)_1\) process are established literature:
\href{https://arxiv.org/abs/math-ph/9907022}{Simon 1999} (the
bounded-below case cited and its theorem passage read in the 1998 test),
and \href{https://math.caltech.edu/simon/biblio.html}{Rosen--Simon 1976}
(bibliographic record). The real-time quartic oscillator also has an
oscillatory-integral treatment by
\href{https://archive.numdam.org/articles/10.1016/j.crma.2005.09.024/}{Albeverio--Mazzucchi
2005} (abstract). None of these references is a novelty claim.

\hypertarget{a-vacuum-and-a-gap-from-the-hamiltonian}{%
\subsection{2. A vacuum and a gap from the
Hamiltonian}\label{a-vacuum-and-a-gap-from-the-hamiltonian}}

Define on \(L^2(\mathbb R,dq)\) the Friedrichs operator
\[ H=-\frac{\hbar^2}{2m}\frac{d^2}{dq^2}+\frac{\lambda}{4}q^4.\tag{2}\]
The form domain is \(H^1(\mathbb R)\cap L^2(\mathbb R,q^4dq)\). The form
is closed and nonnegative. It has compact embedding into \(L^2\):
Rellich compactness holds on every bounded interval, while its unit ball
satisfies
\[\int_{|q|>R}|\psi(q)|^2dq\le R^{-4}\int q^4|\psi(q)|^2dq.\tag{3}\]
Thus \(H\) has compact resolvent. Feynman--Kac makes \(e^{-tH/\hbar}\)
positivity improving for \(t>0\): the free heat kernel is strictly
positive and every finite bridge has a strictly positive potential
weight. By the positivity-improving ground-state theorem, the bottom
eigenvalue \(E_0\) is simple and its normalized eigenfunction \(\Omega\)
is strictly positive. Discreteness implies \(E_1>E_0\) and
\[\Delta=E_1-E_0>0.\tag{4}\]

Writing \(\ell=(\hbar^2/(m\lambda))^{1/6}\) and \(q=\ell x\), unitary
rescaling gives \[H=\hbar^{4/3}\lambda^{1/3}m^{-2/3}
\left(-\tfrac12\frac{d^2}{dx^2}+\tfrac14x^4\right)\tag{5}\] up to that
unitary conjugation. Thus the gap is a positive dimensionless eigenvalue
difference times \(\hbar^{4/3}\lambda^{1/3}m^{-2/3}\); it vanishes in
the indicated small-\(\hbar\) or small-\(\lambda\) limits. This does not
imply positivity of \(\hbar\) from classical premises.

\hypertarget{arbitrary-time-refinement}{%
\subsection{3. Arbitrary time
refinement}\label{arbitrary-time-refinement}}

For \(s>0\), let \[p_s(x,y)=\left(\frac{m}{2\pi\hbar s}\right)^{1/2}
e^{-m(x-y)^2/(2\hbar s)}\] and define the symmetric single-step transfer
kernel \[K_s^{\rm lat}(x,y)=p_s(x,y)
e^{-s[V(x)+V(y)]/(2\hbar)}.\tag{6}\] For a partition \(\pi\) of
\([0,T]\), their product and integration over internal nodes equal
\(p_T(x,y)\) times the expectation, under a Brownian bridge of variance
parameter \(\hbar/m\) from \(x\) to \(y\), of \(\exp[-W_\pi/\hbar]\),
where \[W_\pi=\sum_j\frac{t_{j+1}-t_j}{2}
       [V(q(t_j))+V(q(t_{j+1}))].\tag{7}\] The bridge is continuous
almost surely, so its polygonal interpolation converges uniformly and
\(W_\pi\to\int_0^T V(q(t))dt\) for every sequence with \(|\pi|\to0\).
Since \(0\le e^{-W_\pi/\hbar}\le1\), dominated convergence gives the
Feynman--Kac bridge, including bounded continuous path observables. This
is the argument of the 1998 test, extended from zero to arbitrary fixed
endpoints; no change of \(\hbar\) accompanies the refinement.

To obtain vacuum rather than fixed-endpoint data, weight the two ends by
\(\Omega(x)\Omega(y)\,dx\,dy\) and normalize. The free heat kernel
bounds the tilted kernels, and the \(L^2\) contraction of the free
semigroup gives an integrable bound against \(\Omega(x)\Omega(y)\).
Dominated convergence therefore applies to these finite-interval
distributions. Their continuum normalizer is
\(\langle\Omega,e^{-TH/\hbar}\Omega\rangle=e^{-TE_0/\hbar}\).
Ground-state boundary weights make the continuum restrictions consistent
as the interval grows. Kolmogorov consistency and the locally continuous
bridge laws give a stationary process \(q(\tau)\) on the full time line.
This construction uses the continuum ground state at the boundary; it
does not assert uniform interchange of lattice-spacing and infinite-time
limits for arbitrary lattice boundary states.

Equivalently, with \(A=H-E_0\) and kernel \(K_t\) of \(e^{-tH/\hbar}\),
the stationary one-time law and transition are
\[\mu(dq)=\Omega(q)^2dq,\qquad
 P_t(x,dy)=e^{tE_0/\hbar}\frac{\Omega(y)}{\Omega(x)}K_t(x,y)dy.\tag{8}\]
The eigenvalue equation makes \(P_t\) Markov; symmetry of \(K_t\) makes
it reversible for \(\mu\). Times \(\tau_1\le\cdots\le\tau_n\) and
bounded functions \(f_j\) give
\[\mathbb E\!\left[\prod_{j=1}^n f_j(q(\tau_j))\right]
=\langle\Omega,f_1(Q)e^{-(\tau_2-\tau_1)A/\hbar}f_2(Q)
 \cdots f_n(Q)\Omega\rangle.\tag{9}\]

\hypertarget{reflection-positivity-and-non-gaussianity}{%
\subsection{4. Reflection positivity and
non-Gaussianity}\label{reflection-positivity-and-non-gaussianity}}

Let \(F\) be a finite linear combination of products of bounded
functions of \(q(t)\) at nonnegative times, and let \(\Theta\) reflect
time and complex-conjugate the observable. From (9),
\[\mathbb E[(\Theta F)F]=\|v_F\|^2\ge0,\tag{10}\] where for a single
term with \(0\le t_1\le\cdots\le t_n\),
\[v_F=e^{-t_1A/\hbar}f_1(Q)e^{-(t_2-t_1)A/\hbar}
       \cdots f_n(Q)\Omega,\] and \(v_F\) is extended linearly. This is
reflection positivity for the bounded cylinder algebra. The associated
reconstruction gives the semigroup \(e^{-tA/\hbar}\) and Hilbert space
\(L^2(\mathbb R)\): the vectors \(f(Q)\Omega\), for bounded \(f\), are
dense because \(\Omega>0\) almost everywhere. The coordinate \(Q\) is
self-adjoint as multiplication by \(q\) and \(e^{-itA/\hbar}\) gives
real-time unitary dynamics.

The potential is even, so uniqueness and positivity make \(\Omega\)
even. The field is not Gaussian. If the one-time density \(\Omega^2\)
were Gaussian, positivity would make \(\Omega=C e^{-\alpha q^2/2}\),
\(\alpha>0\). But substituting this into \(H\Omega=E_0\Omega\) leaves
the nonzero quartic term \(\lambda q^4/4\), which cannot be cancelled by
the constant and quadratic terms from the derivative. Thus even the
one-time marginal is non-Gaussian. This is a continuum witness, unlike
the unscaled shrinking cell cumulant in
\href{alpoge-balaban-newton-peek.md}{the Alpoge--Balaban peek}.

For bounded real \(f,g\), spectral projection also gives
\[|\operatorname{Cov}_\mu(f(q(0)),g(q(t)))|
\le e^{-t\Delta/\hbar}\,\|f\Omega\|\,\|g\Omega\|.\tag{11}\] The gap is
an energy gap of this explicitly constructed quantum mechanical
Hamiltonian, not a four-dimensional field-theory mass gap. The claim
here is the bounded-cylinder reconstruction; it does not separately
verify a chosen full hierarchy of unbounded-field growth axioms.

\hypertarget{the-gauge-theory-rung-and-relativity}{%
\subsection{5. The gauge-theory rung and
relativity}\label{the-gauge-theory-rung-and-relativity}}

A literal pure gauge connection on a one-dimensional base has no
two-form curvature: \(F_{00}=0\). A time circle can retain a global
holonomy. The nontrivial model connected to higher-dimensional
Yang--Mills is instead its dimensional reduction: spatial gauge
components survive as adjoint matrices \(X_i(t)\) and the temporal gauge
field imposes Gauss's law. For \(SU(2)\) and \(D\ge2\) matrices, the
\href{low-dimensional-mass-gap.md\#32-the-su2-matrix-model-and-its-abelian-valleys}{matrix
theorem} proves compact resolvent, a simple positive vacuum, and a gap
on the full configuration space and its gauge-invariant subspace. Its
potential is a sum of squared commutators and has unbounded commuting
valleys, so the scalar argument \(V\to\infty\) cannot be reused. The
multicomponent version of (6)--(10), restricted to bounded invariant
functions, supplies the stationary reflection-positive gauge-singlet
construction; positivity and simplicity make the ground state invariant.
This finite-dimensional matrix model keeps only spatially constant
modes; it does not justify discarding the other modes of the full theory
or construct Yang--Mills on \(\mathbb R^4\). No \(SU(3)\) matrix gap is
asserted here.

The field theory (1) is a scalar theory on a one-dimensional spacetime,
and is ordinary nonrelativistic quantum mechanics when \(q\) is read as
the coordinate of a particle. There are no nontrivial boosts in \(0+1\)
spacetime dimensions. A relativistic particle in \(d\) target-space
dimensions instead can use the positive-energy Hamiltonian
\[H_{\rm rel}=\sqrt{c^2P^2+m^2c^4}-mc^2+V(Q).\tag{12}\] Self-adjointness
gives a unitary time evolution when the operator is well-defined, and
the positive semigroup can have a subordinate-Brownian, jump-process
representation under suitable potential assumptions. See
\href{https://arxiv.org/abs/1109.5829}{Hiroshima--Ichinose--Lorinczi}
(abstract and introduction) and
\href{https://arxiv.org/abs/1605.01791}{Gheryan et al.} (abstract),
which treats both relativistic and nonrelativistic Nelson models. Their
theorems have their own assumptions. A confining external \(V\) selects
a frame; it does not establish full Poincare covariance. For \(V=0\) on
\(\mathbb R^d\), the bottom of (12) is not a normalizable vacuum, so the
stationary ground-state construction above cannot be applied without
additional structure. A proper-time worldline and the
Klein--Gordon/Feynman propagator require their own constraint and
normalization conventions.

\hypertarget{consequence-for-state}{%
\subsection{6. Consequence for STATE}\label{consequence-for-state}}

The preliminary \(0+1\) existence statement is a standard theorem for
specified quantum weights and is now assembled with vacuum, reflection
positivity, non-Gaussianity and the continuum limit. The same
architecture connects to the existing \(SU(2)\) matrix-model gap. The
open Newton question remains whether independently justified premises
require a positive universal action scale; the \(\hbar=0\) alternatives
and the real-time/classical distinction remain essential.
\end{document}
